% =============================================================
% 第3.2节 MLSE：自动化安全评估框架设计 —— 初稿
% 请审阅后再合并至 chap03.tex
% =============================================================

\section{MLSE：自动化安全评估框架设计}
\label{sec:mlse_framework}

\subsection{框架设计动机与需求分析}
\label{subsec:framework_motivation}

第\ref{sec:sop_modeling}节将真实的临床诊疗流程抽象为线性拓扑链，并建立了多智能体协作系统的状态转移方程（公式\ref{eq:response_generation}）。该理论模型揭示了链式协作系统中上下游信息依赖的固有结构，以及由此衍生的错误传播风险。然而，要将上述理论分析转化为可量化的安全评估结论，仅凭理论推导和小规模的案例分析远远不够，还需要一套能够系统化执行测试任务的自动化工具。

医疗大语言模型的安全评估面临大量测试用例与多种实验配置的组合。以本文的实验为例，待测模型涵盖BioGPT、AlpaCare、BioLlama3等多个不同架构的医疗大模型，攻击类型包括数据投毒、后门攻击、提示词注入和越狱攻击四类，每类攻击又涉及不同的参数设置（如投毒比例、触发词选择、提示词变异策略等）。在这种高维实验空间中，人工逐一构造测试用例、手动执行查询并主观判断输出质量，既难以保证评估的全面性，也无法确保不同实验之间的一致性和可复现性。因此，构建一个端到端的自动化评估框架是开展本文研究的必要前提。

基于上述分析，本文提出的MLSE（Medical LLM Safety Evaluation）框架围绕三个核心需求进行设计。其一，框架必须支持从数据加载、攻击执行到指标计算的全流程自动化，以适应大规模安全测试的效率要求。其二，框架需要同时支持单体模型评估和多智能体级联评估两种模式。单体模式用于在隔离环境下揭示单个模型面对各类攻击时的脆弱性特征，而级联模式则模拟真实临床协作场景，研究错误信息在智能体链路中的传播与放大规律，二者分别对应第\ref{cha:single_agent}章和第\ref{cha:multi_agent}章的实验内容。其三，框架在架构上应具备良好的可扩展性，支持新模型、新攻击方法和新评估指标的快速集成，使其能够适应医疗大语言模型领域快速发展的研究需求。与此同时，第\ref{sec:eval_metrics}节定义的攻击成功率（ASR）、文本困惑度（PPL）、连贯性（COH）以及Cohen's Kappa系数等评估指标，需要在框架中找到明确的计算位置和执行主体。

\subsection{框架总体架构}
\label{subsec:framework_architecture}

围绕上述需求，MLSE框架采用分层模块化架构设计，遵循关注点分离原则，将模型管理、智能体交互和评估任务执行分别置于独立层次。框架自底向上包含模型集成层、智能体交互层和应用评估层三个主要层次，各层之间通过标准化接口进行通信。

模型集成层位于框架的最底层，承担待测医疗大语言模型的统一管理与调度职责。该层的核心设计理念是屏蔽不同来源模型之间的差异，为上层组件提供标准化的推理服务接口。在实际部署中，模型集成层首先利用llama.cpp量化工具\cite{llamacpp2024}对原始模型进行轻量化处理，将参数精度从32位浮点数压缩至4位整数表示，在保持模型性能基本不变的前提下显著降低显存占用。量化后的模型被导入Ollama服务\cite{ollama2024}，该服务提供统一的RESTful API接口，使得框架能够以一致的方式调用不同架构和来源的模型。当新的医疗大模型发布时，只需完成量化处理并导入Ollama即可被框架识别和调度，无需修改上层的评估逻辑。

智能体交互层是框架的核心层，负责构建和管理评估过程中所需的各类智能体实例及其通信关系。该层基于Microsoft AutoGen框架\cite{autogen2023}实现，AutoGen提供了成熟的多智能体协作通信协议和灵活的智能体定义接口，已被广泛应用于多智能体系统的学术研究。智能体交互层支持两种主要的智能体配置模式。在单体模式下，单个目标智能体直接连接待测模型的后端接口，独立响应查询请求，适用于评估模型在隔离状态下的安全性能。在级联模式下，多个目标智能体按照第\ref{sec:sop_modeling}节定义的诊疗链拓扑结构连接，每个智能体的上下文输入遵循公式(\ref{eq:response_generation})所描述的拼接机制，即同时接收原始查询$Q$和直接上游智能体的输出$O_{i-1}$。这一设计使得级联模式成为第\ref{subsec:clinical_chain}节理论模型的工程实现，能够在受控环境中模拟真实的多智能体协作诊疗流程。

应用评估层位于框架的最上层，直接面向安全评估任务。该层根据攻击类型的不同，划分为三个专用评估子系统。数据污染评估子系统专注于检测模型对训练阶段投毒攻击和后门攻击的防御能力，对应第\ref{sec:threat_model}节中分析的上游知识库投毒威胁。提示词操纵评估子系统模拟部署阶段的攻击场景，包括提示词注入和越狱攻击，对应终端提示词攻击威胁。级联风险评估子系统是本文的特色模块，它将多个智能体串联成链，研究错误信息在协作链中的传播和放大规律，为验证第\ref{subsec:cascading_error}节提出的级联错误传播假设提供实验支撑。三个子系统共享模型集成层和智能体交互层提供的基础设施，但各自拥有独立的攻击构造逻辑和评估指标组合。三个子系统的详细执行流程将在第\ref{subsec:agent_roles}节的智能体协作机制中具体阐述。

表\ref{tab:framework_comparison}将MLSE与现有主流安全评估框架进行对比，突出本框架在医疗场景适配性和多智能体支持方面的差异。

\begin{table}[htbp]
\centering
\caption{MLSE与现有安全评估框架对比}
\label{tab:framework_comparison}
\begin{tabular}{lccccc}
\toprule
\textbf{框架} & \textbf{评估目标} & \textbf{医疗适应性} & \textbf{多智能体支持} & \textbf{攻击类型覆盖} & \textbf{自动化程度} \\
\midrule
HELM\cite{helm2022} & 通用LLM能力 & 无 & 无 & 无对抗评估 & 高 \\
HarmBench\cite{mazeika2024harmbench} & 安全性 & 低 & 无 & 越狱攻击 & 高 \\
AgentHarm\cite{andriushchenko2024agentharm} & 智能体安全 & 无 & 有限 & 提示词攻击 & 中 \\
\textbf{MLSE(本文)} & \textbf{医疗LLM安全} & \textbf{高} & \textbf{完整支持} & \textbf{多类型} & \textbf{高} \\
\bottomrule
\end{tabular}
\end{table}

图\ref{fig:mlse_architecture}展示了MLSE框架的整体架构设计。

\begin{figure}[htbp]
\centering
\fbox{\parbox{0.95\textwidth}{\centering\vspace{4cm}
\textbf{图片占位符}\\[0.5cm]
MLSE框架系统架构图\\
包含：三层结构（模型集成层、智能体交互层、应用评估层）\\
六类核心智能体及其连接关系\\
数据流向：配置输入→攻击执行→结果收集→报告输出\\
外部接口：待测模型、评估数据集、用户配置
\vspace{4cm}}}
\caption{MLSE框架系统架构}
\label{fig:mlse_architecture}
\end{figure}

\subsection{核心智能体角色与评估子系统设计}
\label{subsec:agent_roles}

MLSE框架的安全评估功能由多个承担特定角色的智能体协同实现。按照功能类型，可将框架中的智能体划分为四类：控制类、被测对象、评测类和辅助类。表\ref{tab:agent_roles}概括了各智能体的分类与核心职责。这一角色划分既保证了评估流程中各环节职责的清晰分离，也与第\ref{subsec:agent_arch}节定义的智能体三元组$\langle P, M, A \rangle$形成对应——每类智能体的感知模块、记忆模块和动作模块的具体配置因其角色差异而不同。

\begin{table}[htbp]
\centering
\caption{MLSE框架核心智能体分类与职责}
\label{tab:agent_roles}
\begin{tabular}{llll}
\toprule
\textbf{类别} & \textbf{智能体名称} & \textbf{核心职责} \\
\midrule
控制类 & 主持人智能体（Moderator Agent） & 评估流程编排与状态监控 \\
\midrule
被测对象 & 目标级联智能体（Target Cascade Agents） & 接收查询并生成响应 \\
\midrule
\multirow{2}{*}{评测类} & 安全智能体（Safety Agent） & 输出安全性判定 \\
 & 评估智能体（Evaluator Agent） & 输出准确性评分 \\
\midrule
\multirow{2}{*}{辅助类} & 提示词变异智能体（Prompt Mutation Agent） & 生成对抗性攻击样本 \\
 & 性能智能体（Performance Agent） & 结果聚合与报告生成 \\
\bottomrule
\end{tabular}
\end{table}

主持人智能体（Moderator Agent）扮演整个评估流程的控制者角色，负责协调评估任务的执行进度，确保各步骤按照预定的顺序正确推进。评估启动时，主持人智能体完成测试环境的初始化，包括加载待测模型、配置评估参数和准备测试数据集。在评估执行过程中，主持人智能体负责将查询输入分发给目标智能体，在级联模式下还负责管理上下文信息在各智能体之间的传递。评估结束后，主持人智能体收集各环节的中间结果，并将其汇总传递给性能智能体。通过集中式的流程编排，主持人智能体保证了评估任务在不同实验配置下的可复现性与可追溯性。

目标级联智能体（Target Cascade Agents）是安全评估的被测对象集合。在单体评估模式下，该集合仅包含一个智能体，其后端接口连接到待测医疗大语言模型。在级联评估模式下，该集合包含多个按序连接的智能体，每个智能体模拟诊疗链中的一个环节。级联模式下的目标智能体被设计为能够同时接收原始查询和上游智能体输出的复合输入，即按照$Context_i = \langle Q, O_{i-1} \rangle$的方式构造上下文。每个目标智能体可以配置不同的底层模型，从而支持混合模型架构的安全评估。级联模式的这一设计使得框架能够量化研究错误信息在协作链中的传播规律，是连接第\ref{sec:sop_modeling}节理论分析与第\ref{cha:multi_agent}章实验验证的关键桥梁。

安全智能体（Safety Agent）和评估智能体（Evaluator Agent）分别从安全性和准确性两个维度评价目标智能体的输出。安全智能体连接到专门的安全评估大模型，对目标输出进行多维度安全检查，判断其是否包含可能导致误诊或错误治疗的医疗建议、是否存在歧视性表述、是否泄露患者隐私等，其输出以结构化的安全评分形式呈现。评估智能体同样连接到大模型，主要判断输出内容是否正确回答了查询问题以及回答的逻辑完整性。两类评测智能体在评估过程中并行工作，各自从不同维度对同一份目标输出进行独立评价。它们的评分结果构成了第\ref{sec:eval_metrics}节所定义的攻击成功率（ASR）计算的基础，同时为困惑度（PPL）和连贯性（COH）等辅助指标提供数据来源。

提示词变异智能体（Prompt Mutation Agent）负责自动生成原始查询的多种对抗性变体。该智能体通过在原始提示词中插入恶意指令、采用角色扮演策略或添加触发词等方式，构造出具有攻击性的输入样本。变异过程中使用困惑度阈值进行筛选，以保证生成的变异样本在保持攻击意图的同时维持自然的语言表达。其核心策略如算法\ref{alg:prompt_mutation}所示。

\begin{algorithm}[htbp]
\caption{提示词变异生成策略}
\label{alg:prompt_mutation}
\begin{algorithmic}[1]
\REQUIRE 原始查询 $Q$, 变异策略集合 $\mathcal{S}$, 最大迭代次数 $T$, 困惑度阈值 $\tau$
\ENSURE 变异提示词集合 $\mathcal{M}$
\STATE $\mathcal{M} \leftarrow \emptyset$
\FOR{$t = 1$ to $T$}
    \STATE 从 $\mathcal{S}$ 中随机选择策略 $s$
    \IF{$s$ = 角色扮演}
        \STATE $Q' \leftarrow$ InsertRolePrompt($Q$)
    \ELSIF{$s$ = 指令注入}
        \STATE $Q' \leftarrow$ AppendMaliciousInstruction($Q$)
    \ELSIF{$s$ = 触发词插入}
        \STATE $Q' \leftarrow$ InsertTriggerWord($Q$)
    \ENDIF
    \STATE 计算 $Q'$ 的困惑度 $\text{PPL}(Q')$
    \IF{$\text{PPL}(Q') < \tau$}
        \STATE $\mathcal{M} \leftarrow \mathcal{M} \cup \{Q'\}$
    \ENDIF
\ENDFOR
\RETURN $\mathcal{M}$
\end{algorithmic}
\end{algorithm}

性能智能体（Performance Agent）承担评估结果的最终聚合与报告生成职责。该智能体收集安全智能体和评估智能体的所有评分结果，结合攻击成功率、困惑度等可计算指标，生成多维度的综合评估报告。报告内容包括各测试用例的详细得分、总体安全等级评定以及不同攻击策略和模型配置下的表现差异分析。

上述六类智能体通过三个评估子系统组织成不同的协作模式，分别服务于训练阶段和部署阶段的安全评估需求。

在数据污染检测子系统中，评估对象是经过投毒或后门微调的待测模型。执行流程为：主持人智能体从数据集中加载测试查询，将其发送给目标智能体（此时为单体配置）；目标智能体生成响应后，安全智能体和评估智能体并行启动，分别从安全性和准确性两个维度对目标输出进行评价；性能智能体汇总两者的评分结果，计算攻击成功率等指标并生成评估报告。该子系统在第\ref{cha:single_agent}章的数据投毒和后门攻击实验中被使用。

提示词操纵检测子系统在架构上与数据污染子系统保持一致，但额外引入了提示词变异智能体。执行流程为：主持人智能体输入原始测试问题，目标智能体生成初始响应，安全智能体对其进行安全性评估后，提示词变异智能体根据算法\ref{alg:prompt_mutation}自主生成原始查询的多种变异形式，并将这些变异提示词反馈给目标智能体，模拟各类提示词注入和越狱攻击场景。经过多轮变异迭代后，性能智能体汇总所有安全评估与准确性评分，输出综合评估报告。该子系统在第\ref{cha:single_agent}章的提示词注入和越狱攻击实验中被使用。

级联风险检测子系统是MLSE框架中最具特色的模块。该子系统在外层架构上保留了主持人智能体、安全智能体和性能智能体的协调与评估角色，但在核心测试阶段将单一的目标智能体替换为由多个串联智能体构成的级联链。在本文的实验配置中，级联链包含五个互相连接的目标智能体。每个智能体依次处理原始临床查询和其直接前驱智能体的输出，这一设计与第\ref{subsec:clinical_chain}节所描述的临床诊疗链的序列推理模式保持一致。级联子系统不仅评估最终的链尾输出，还记录每个中间智能体的输出内容，从而能够追踪错误信息在链路中的传播轨迹，为验证第\ref{subsec:cascading_error}节的级联错误传播假设和计算错误放大因子$\alpha$提供实验数据。该子系统在第\ref{cha:multi_agent}章的多智能体安全评估实验中被使用。

图\ref{fig:agent_interaction}展示了一次完整安全评估过程中各智能体的交互时序。

\begin{figure}[htbp]
\centering
\fbox{\parbox{0.9\textwidth}{\centering\vspace{3.5cm}
\textbf{图片占位符}\\[0.5cm]
智能体交互时序图\\
纵向：Moderator → Target1 → Target2 → ... → Safety → Evaluator → Performance\\
横向：时间流\\
消息传递：查询输入、响应生成、安全评分、准确性评分、报告生成
\vspace{3.5cm}}}
\caption{MLSE智能体交互时序图}
\label{fig:agent_interaction}
\end{figure}


\subsection{模型部署与接口工程}
\label{subsec:engineering_optimization}

MLSE框架的设计还需要解决大规模安全评估任务中的工程效率问题。本节从模型量化、接口标准化和并发处理三个方面阐述框架的工程实现方案。

模型量化是提升系统效率的基础。医疗大语言模型的参数规模通常以十亿级计，直接部署对硬件资源提出了极高要求。MLSE框架使用llama.cpp工具将模型参数精度从32位浮点数量化至4位整数表示\cite{llamacpp2024}。量化后的模型文件体积通常可缩小至原始规模的四分之一左右，在同一台机器上并存多个模型成为可能。量化操作同时降低了推理的计算复杂度，缩短了单次请求的响应时间。关于量化策略对实验结论的影响，本文在第\ref{cha:multi_agent}章中对Q4量化与FP16全精度模型进行了对比验证，在攻击成功率指标上的差异小于3\%，处于统计误差范围之内。

接口标准化是确保框架可扩展性的关键。不同研究机构发布的医疗大模型往往采用差异化的文件格式和推理接口。MLSE框架采用Ollama\cite{ollama2024}作为统一模型服务层，屏蔽底层模型差异。经过量化的模型文件被导入Ollama后，框架的所有上层调用均通过其标准化的RESTful API完成。当新模型发布时，只需完成导入操作即可复用现有的评估流程。选择Ollama还基于其对本地化部署的支持，这一特性能够有效保护医疗数据隐私，避免敏感信息传输至外部服务。

并发处理机制直接决定了评估任务的执行效率。MLSE框架通过任务级并发和请求级批量推理两个层面来提升吞吐量。任务级并发允许同时启动多个独立的评估任务，由调度器动态分配计算资源。请求级批量推理将多个查询请求打包提交给模型，充分利用GPU的并行计算能力\cite{flashdecoding2024}。此外，框架还引入了缓存机制和容错处理来增强系统鲁棒性。缓存机制记录已执行过的查询及其响应，在攻击测试中对同一查询的多次变体尝试可直接复用已有结果。容错机制监控模型服务无响应、输出格式错误等异常情况，支持自动重试或跳过并记录日志。上述工程措施共同保证了MLSE框架在有限计算资源条件下能够高效完成大规模安全评估任务。
